\scanpage{63}
% block: numerical data continued
et la classe de $\operatorname{End}(\deleted{\unsure{S|B}}\handwritten{M})$ dans $\operatorname{Br}(Z_{\handwritten{\sigma}})$ (le couple $(Z_{\handwritten{\sigma}},\handwritten{\xi_\sigma})$ étant défini à isomorphisme près). Sans supposer $\handwritten{\mathcal{M}}$ semi\handwritten{\unsure{(}}-si\deleted{p}mple, on peut\inserted{\handwritten{encore\begin{tikzpicture}[x=1cm,y=1cm,baseline=-.1cm]\draw[thin] (0,.15).. controls (.1,-.1) and (-.2,-.1)..(-1,-.1)--(-2.1,-.1).. controls (-2.3,-.1) and (-2.35,-.2)..(-2.35,-.5);\end{tikzpicture}}} dire que la plus grande sous-catégorie semi-simple(engend\deleted{g}rée par les objets semi-simples) est connue à équiva\unsure{\deleted{e}l|\deleted{l}e}ence près par la connaissance précédente, qui implique donc la connaissance de $K(\handwritten{\mathcal{M}})=\handwritten{\mathbb{Z}^{(\Sigma)}}$,\inserted{\handwritten{sa partie effective,\begin{tikzpicture}[x=1cm,y=1cm,baseline=-.1cm]\draw[thin] (0,.15).. controls (.1,-.1) and (-.2,-.1)..(-1,-.1)--(-2.1,-.1).. controls (-2.3,-.1) and (-2.35,-.2)..(-2.35,-.5);\end{tikzpicture}}} et \unsure{\deleted{d}s|\deleted{s}d}a variance pour \unsure{\deleted{d}l|\deleted{l}d}es changements de\deleted{s} base par $k'$ variable sur $k$. De façon précise,$\handwritten{\Sigma'=\Sigma(\mathcal{M}_{k'})}$ en tant qu'ensemble sur $\handwritten{\Sigma}$ correspond à l'ensemble \deleted{des corps} dont la fibre \deleted{est} en $\handwritten{\sigma}$ est l'ensemble des corps composants de $Z_{\handwritten{\sigma}}\otimes_k k'$, soient $Z'^{i}_{\handwritten{\sigma}}\deleted{\unsure{i|\times}}$, l'application $K(\handwritten{\mathcal{M}})=\underline{Z}^{(\handwritten{\Sigma})}\handwritten{\longrightarrow}K(\handwritten{\mathcal{M}_{k'}})=\underline{Z}^{(\handwritten{\Sigma'})}$ est donné par
\[
\handwritten{\sigma\longmapsto}\ n_i\ (\handwritten{\sigma},i)\ ,
\]
où $n_i$ est donné par la règle suivante: si $d$ est \deleted{le degré de} l'ordre de $\handwritten{\xi_\sigma\in}\operatorname{Br}(Z_{\handwritten{\sigma}})$, celui de son image dans $\operatorname{Br}(Z'^{i}_{\handwritten{\sigma}})$ est un \deleted{idivs} diviseur $d'_i$ de $d$, et on a $n_i=d/d_i$. (A généraliser au cas où les $Z_{\handwritten{\sigma}}$ ne sont pas nécessairement séparables). Si $k'$ est une \deleted{\unsure{clô}} extension alg. close de $k$, l'ensemble $\handwritten{\Sigma'}$ s'identifie à l'ensemble somme des \uline{fibres} en $k'$ des $\operatorname{Spec}(Z_{\handwritten{\sigma}})$ ....

% block: b numerical characters
b) Supposons maintenant $\handwritten{\mathcal{M}}$ \deleted{semi-tannakienne} munie d'un produit tensoriel à contraintes d'associativité et commutativité, avec 1.

On désigne sous le nom de ``caractères numériques\deleted{''} de $\handwritten{\Sigma}$ '' les données suivantes:
\margin{\handwritten{préciser cette notion, en faisant d'un ``caractère numérique'' une structure susceptible de variation fonctorielle....}}

% block: numerical characters item 1
\unsure{a|1}) l'ensemble (à bijection près)\inserted{\handwritten{$\deleted{\unsure{\xi|\Sigma}}\Sigma$\begin{tikzpicture}[x=1cm,y=1cm,baseline=-.1cm]\draw[thin] (-.1,-.25).. controls (0,0) and (.15,.1)..(.3,.4);\end{tikzpicture}}} des \deleted{com} classes d'isomorphie d'objets simples de $\handwritten{\mathcal{M}}$.

% block: numerical characters item 2
2) La donnée de la multiplication de $K(\handwritten{\mathcal{M}})={\unsure{\mathbb{Z}|B|Z}}^{(\handwritten{\Sigma})}$; ce qui revient à la donnée pour tout $\handwritten{\sigma,\tau\in\Sigma}$ du produit
\[
\handwritten{\sigma\tau}=\handwritten{\sum_\rho}c_{\handwritten{\sigma\tau}}^{\handwritten{\rho}}\handwritten{\rho}\qquad\handwritten{(c_{\sigma,\tau}^{\rho}\in\mathbb{N})},
\]
\handwritten{i.e. du système d'entiers naturels $\sigma_{\sigma,\tau}^{\rho}$.}
% archive sheet number: 62
